%% RESUMO EM LINGUA ESTRANGEIRA (elemento obrigatorio -- NBR 6028)
%% Traducao fiel do resumo em portugues. As keywords sao definidas em
%% config/dados.tex.

\begin{abstract}
This dissertation investigates the use of artificial intelligence for
continuous monitoring and compliance with the Brazilian General Data Protection
Law (LGPD) in the Electronic Information System (SEI). It starts from the large
volume of proceedings handled by the system and from the limitations of manual
control, which scales poorly and is prone to failure. The research problem is
how artificial intelligence and natural language processing solutions can
automate the identification of personal and sensitive data in documents. The
criteria for using these techniques are validated by experts through the DELPHI
Method. The research thus seeks to contribute to improving the security and efficiency of
public records management, through a flow that goes from reading SEI documents
to issuing compliance alerts.
\end{abstract}
